% ============================================================
%  Préambule — packages et configuration
% ============================================================

% Encodage et langue
\usepackage[utf8]{inputenc}
\usepackage[T1]{fontenc}
\usepackage[french]{babel}
\usepackage{csquotes}

% Typographie
\usepackage{lmodern}
\usepackage{microtype}
\usepackage[a4paper,top=2.5cm,bottom=2.5cm,inner=2.8cm,outer=2.2cm]{geometry}
\usepackage{setspace}
\onehalfspacing
\setlength{\headheight}{14pt}

% Mathématiques
\usepackage{amsmath,amssymb,amsthm}
\usepackage{mathtools}

% Figures et tableaux
\usepackage{graphicx}
\graphicspath{{figures/}}
\usepackage{booktabs}
\usepackage{tabularx}
\usepackage{multirow}
\usepackage{longtable}
\usepackage{caption}
\usepackage{subcaption}
\captionsetup[figure]{labelfont=bf,format=plain,font=small}
\captionsetup[table]{labelfont=bf,format=plain,font=small}

% Couleurs
\usepackage[dvipsnames,table]{xcolor}
\definecolor{anssi}{HTML}{003366}
\definecolor{soft}{HTML}{2A5D8F}
\definecolor{lightbg}{HTML}{F0F4F8}
\definecolor{codebg}{HTML}{F6F8FA}
\definecolor{codecomment}{HTML}{6A737D}
\definecolor{codestring}{HTML}{032F62}
\definecolor{codekeyword}{HTML}{D73A49}

% Code (listings)
\usepackage{listings}
\lstset{
  basicstyle=\ttfamily\footnotesize,
  backgroundcolor=\color{codebg},
  commentstyle=\color{codecomment}\itshape,
  keywordstyle=\color{codekeyword}\bfseries,
  stringstyle=\color{codestring},
  numbers=left,
  numberstyle=\tiny\color{gray},
  numbersep=8pt,
  frame=single,
  framerule=0pt,
  rulecolor=\color{lightbg},
  breaklines=true,
  showstringspaces=false,
  tabsize=2,
  captionpos=b,
  xleftmargin=10pt,
  inputencoding=utf8,
  extendedchars=true,
  literate=%
    {é}{{\'e}}1 {è}{{\`e}}1 {ê}{{\^e}}1 {ë}{{\"e}}1
    {à}{{\`a}}1 {â}{{\^a}}1 {ä}{{\"a}}1
    {ô}{{\^o}}1 {ö}{{\"o}}1
    {î}{{\^\i}}1 {ï}{{\"\i}}1
    {ù}{{\`u}}1 {û}{{\^u}}1 {ü}{{\"u}}1
    {ç}{{\c{c}}}1 {É}{{\'E}}1 {È}{{\`E}}1
    {œ}{{\oe}}1 {Œ}{{\OE}}1 {æ}{{\ae}}1 {Æ}{{\AE}}1
    {€}{{\euro}}1 {—}{{---}}1 {–}{{--}}1
    {«}{{\guillemotleft{}}}1 {»}{{\guillemotright{}}}1
    {“}{{``}}1 {”}{{''}}1 {‘}{{`}}1 {’}{{'}}1
    {…}{{\ldots}}1 {·}{{$\cdot$}}1
}
\lstdefinelanguage{Python3}{
  language=Python,
  morekeywords={async,await,True,False,None,match,case},
}

% Titres de chapitres / sections
\usepackage{titlesec}
\titleformat{\chapter}[display]
  {\normalfont\Huge\bfseries\color{anssi}}
  {\filright\Large CHAPITRE \thechapter}
  {1ex}
  {\titlerule\vspace{1ex}\filright}
\titleformat{\section}{\normalfont\Large\bfseries\color{anssi}}{\thesection}{1em}{}
\titleformat{\subsection}{\normalfont\large\bfseries\color{soft}}{\thesubsection}{1em}{}
\titleformat{\subsubsection}{\normalfont\normalsize\bfseries\color{soft}}{\thesubsubsection}{1em}{}

% En-têtes et pieds de page
\usepackage{fancyhdr}
\pagestyle{fancy}
\fancyhf{}
\fancyhead[LE]{\nouppercase{\leftmark}}
\fancyhead[RO]{\nouppercase{\rightmark}}
\fancyfoot[LE,RO]{\thepage}
\fancyfoot[RE,LO]{\small\textit{Mémoire PFE — UrbanUpC SOC + HOMO-CI}}
\renewcommand{\headrulewidth}{0.4pt}
\renewcommand{\footrulewidth}{0.2pt}

% Hyperliens
\usepackage{hyperref}
\hypersetup{
  colorlinks=true,
  linkcolor=anssi,
  urlcolor=soft,
  citecolor=anssi,
  pdftitle={\memTitle},
  pdfauthor={\memAuthor},
}
% cleveref désactivé : il plantait fatalement sur les références encore
% indéfinies au premier passage. \autoref (hyperref) fait l'équivalent
% (« chapitre N », « section N ») et n'émet qu'un warning si la ref manque.
% On gère aussi la forme avec liste : \cref{a,b,c} -> "chapitre N, N et N".
\makeatletter
\newcommand{\cref@list}[1]{%
  \@cref@list@helper#1,\@nil%
}
\def\@cref@list@helper#1,#2\@nil{%
  \autoref{#1}%
  \def\@cref@rest{#2}%
  \ifx\@cref@rest\@empty\else
    \@cref@list@continue#2\@nil
  \fi
}
\def\@cref@list@continue#1,#2\@nil{%
  \def\@cref@rest{#2}%
  \ifx\@cref@rest\@empty
    \space et~\autoref{#1}%
  \else
    , \autoref{#1}%
    \@cref@list@continue#2\@nil
  \fi
}
\providecommand{\cref}[1]{\cref@list{#1}}
\providecommand{\Cref}[1]{\cref@list{#1}}
\makeatother
% Noms français pour \autoref
\addto\extrasfrench{%
  \def\chapterautorefname{chapitre}%
  \def\sectionautorefname{section}%
  \def\subsectionautorefname{section}%
  \def\subsubsectionautorefname{section}%
  \def\figureautorefname{figure}%
  \def\tableautorefname{tableau}%
  \def\equationautorefname{équation}%
  \def\partautorefname{partie}%
}

% Boîtes (encadrés) — chargement allégé (pas besoin de [most])
\usepackage[skins,breakable]{tcolorbox}
\newtcolorbox{contributionbox}[1]{
  colback=lightbg,colframe=anssi,fonttitle=\bfseries,
  title=#1,enhanced,boxrule=0.7pt,
}
\newtcolorbox{takeaway}{
  colback=anssi!5,colframe=anssi,fonttitle=\bfseries,
  title=À retenir,enhanced,boxrule=0.7pt,
}
\newtcolorbox{warning}{
  colback=red!5,colframe=red!50!black,fonttitle=\bfseries,
  title=Limite,enhanced,boxrule=0.7pt,
}

% Glossaire / acronymes
\usepackage{acronym}

% Symbols pour les croix/checks
\usepackage{eurosym}
\usepackage{pifont}
\newcommand{\cmark}{\ding{51}}
\newcommand{\xmark}{\ding{55}}

% Espaces verticaux
\setlength{\parskip}{0.6em}
\setlength{\parindent}{0pt}

% Référence de hyperref aux contributions
\newcommand{\contribref}[1]{\textbf{C-#1}}
\newcommand{\hypref}[1]{\textbf{H#1}}
\newcommand{\sqref}[1]{\textbf{SQ-#1}}
\newcommand{\todofig}[1]{\textcolor{red}{[\textbf{FIGURE À AJOUTER :} #1]}}
\newcommand{\todocapture}[1]{\textcolor{red}{[\textbf{CAPTURE À AJOUTER :} #1]}}
